% META: {"id": "TSPE-CONTIN-CO-003", "chapitre": "TSPE-CONTINUITE", "type_objet": "corrige", "exercice_ref": "TSPE-CONTIN-EX-003", "capacites_codes": ["C1"], "sources_inspiration": [], "status": "approved"}
% BEGIN-VERIFY
% from sympy import symbols, diff
% x = symbols('x')
% f = x**3 - 3*x + 3
% assert diff(f, x) == 3*x**2 - 3
% assert f.subs(x, -3) == -15
% assert f.subs(x, -1) == 5
% assert -15 <= 1 <= 5
% # la derivee s'annule en -1 : stricte positivite sur l'ouvert seulement
% assert diff(f, x).subs(x, -1) == 0
% assert all(diff(f, x).subs(x, v) > 0 for v in (-3, -2, -1.5, -1.01))
% END-VERIFY
\begin{corrige}{TSPE-CONTIN-EX-003}

\textbf{Q1.} $f'(x) = 3x^2 - 3 = 3(x-1)(x+1)$. Pour $x \in \left]-3,-1\right[$, on a $-3 < x < -1$, donc $|x| > 1$, donc $x^2 > 1$ et $f'(x) = 3(x^2-1) > 0$. En revanche $f'(-1) = 3 - 3 = 0$ : la derivee s'annule a la borne droite.

\textbf{Q2.} $f$ est continue sur $[-3,-1]$ et sa derivee est strictement positive sur l'intervalle ouvert $\left]-3,-1\right[$. D'apres le theoreme reliant signe de la derivee et sens de variation, $f$ est strictement croissante sur $[-3,-1]$ tout entier : l'annulation de $f'$ en la seule borne $-1$ n'empeche pas la stricte croissance sur l'intervalle ferme.

\textbf{Q3.} $f(-3) = -27 + 9 + 3 = -15$. $f(-1) = -1 + 3 + 3 = 5$.

\textbf{Q4.} $f$ est continue et strictement croissante sur $[-3,-1]$, et $1$ est compris entre $f(-3)=-15$ et $f(-1)=5$. Par le TVI (cas strictement monotone), l'equation $f(x)=1$ admet une unique solution $\alpha \in [-3,-1]$.

\end{corrige}
